\documentclass[12pt]{article}
\usepackage[utf8]{inputenc}
\usepackage[T1]{fontenc}
\usepackage{tikz}
\usepackage{xcolor}

\usetikzlibrary{positioning, calc}
\pagestyle{empty}

% Farbdefinitionen nach deiner Vorlage
\definecolor{formyellow}{HTML}{FFF9C4} % Authentisches Formular-Gelb
\definecolor{formred}{HTML}{E57373}   % Helles Rot für Rahmen und feste Label
\definecolor{formgreen}{HTML}{2E7D32} % Grün für den ausfüllbaren Text

\begin{document}
	
	\begin{tikzpicture}[remember picture, overlay, shift=(current page.south west)]
		
		% 1. Gelber Hintergrund
		\fill[formyellow] (1.5cm, 1.5cm) rectangle (\paperwidth-1.5cm, \paperheight-1.5cm);
		
		% =========================================================================
		% FÄLSCHUNGSSCHUTZ / SATIRE-WASSERZEICHEN (Dezent im Hintergrund)
		% =========================================================================
		\node[rotate=35, scale=4.5, opacity=0.05, text=black, font=\sffamily\bfseries] 
		at (\paperwidth/2, \paperheight/2 + 2cm) {EINLADUNG -- KEIN DOKUMENT};
		\node[rotate=35, scale=4.5, opacity=0.05, text=black, font=\sffamily\bfseries] 
		at (\paperwidth/2, \paperheight/2 - 3cm) {SATIRE -- PARTY-GAG};
		
		% 2. Vertikaler grüner Text ganz links
		\node[rotate=90, font=\sffamily\bfseries\small, color=formgreen, anchor=south] 
		at (2.1cm, \paperheight/2) {Bei nicht rechtzeitiger Vorlage droht extreme Langeweile!};
		
		% =========================================================================
		% RECHTE SEITE: TITEL & HINWEISE
		% =========================================================================
		
		% Haupttitel
		\node[anchor=north west, font=\sffamily\bfseries\Large, color=formgreen, align=left] 
		at (9.5cm, \paperheight-1.8cm) {Feier-\\bescheinigung};
		
		\node[anchor=north east, font=\sffamily\bfseries\Large] 
		at (\paperwidth-2cm, \paperheight-1.8cm) {1};
		
		% Box: Ausfertigung zur Vorlage...
		\draw[formred, thick] (9.5cm, \paperheight-3.8cm) rectangle (\paperwidth-2cm, \paperheight-4.8cm);
		\node[anchor=center, font=\sffamily\bfseries\small, color=formgreen, align=center] 
		at ($(9.5cm, \paperheight-3.8cm)!0.5!(\paperwidth-2cm, \paperheight-4.8cm)$) 
		{Ausfertigung zur Vorlage\\beim Gastgeber};
		
		% Gedicht rechts
		\node[anchor=north west, font=\sffamily\small, color=formgreen, align=left, text width=9cm] 
		at (9.5cm, \paperheight-5.2cm) {
			Wer sich fragt: „Was soll ich kaufen?“\\
			Braucht sich nicht die Haare raufen.\\
			Willst du mir eine Freude machen,\\
			lass doch mein Sparschwein lachen.\\[1em]
			\footnotesize Bitte um Rückantwort - bis 01.12.21\\
			Wer nicht kann, ruft an - 0111/123456
		};
		
		% =========================================================================
		% LINKE SEITE: FORMULARFELDER
		% =========================================================================
		
		% Box 1: Anlass
		\draw[formred, thick] (2.7cm, \paperheight-1.8cm) rectangle (9.0cm, \paperheight-3.0cm);
		\node[anchor=north west, font=\sffamily\tiny\bfseries, color=formred] at (2.8cm, \paperheight-1.9cm) {Anlass};
		\node[anchor=west, font=\sffamily\small, color=formgreen] at (2.8cm, \paperheight-2.6cm) {Einladung zum 30.ten Geburtstag};
		
		% Box 2: Gastgeber & Geb. am
		\draw[formred, thick] (2.7cm, \paperheight-3.0cm) rectangle (9.0cm, \paperheight-4.2cm);
		\draw[formred, thick] (7.0cm, \paperheight-3.0cm) -- (7.0cm, \paperheight-4.2cm);
		
		\node[anchor=north west, font=\sffamily\tiny\bfseries, color=formred] at (2.8cm, \paperheight-3.1cm) {Gastgeber};
		\node[anchor=west, font=\sffamily\small, color=formgreen] at (2.8cm, \paperheight-3.8cm) {Max Mustermann};
		
		\node[anchor=north west, font=\sffamily\tiny\bfseries, color=formred] at (7.1cm, \paperheight-3.1cm) {geb. am};
		\node[anchor=west, font=\sffamily\small, color=formgreen] at (7.1cm, \paperheight-3.8cm) {11.12.1991};
		
		% Box 3: Veranstaltungsort
		\draw[formred, thick] (2.7cm, \paperheight-4.2cm) rectangle (9.0cm, \paperheight-6.4cm);
		\node[anchor=north west, font=\sffamily\tiny\bfseries, color=formred] at (2.8cm, \paperheight-4.3cm) {Veranstaltungsort};
		\node[anchor=north west, font=\sffamily\small, color=formgreen, align=left] at (2.8cm, \paperheight-4.8cm) {
			Auf der Tulpe\\
			Tulpenweg 60\\
			90768 Fürth
		};
		
		% Box 4: Motto
		\draw[formred, thick] (2.7cm, \paperheight-6.4cm) rectangle (9.0cm, \paperheight-7.6cm);
		\node[anchor=north west, font=\sffamily\tiny\bfseries, color=formred] at (2.8cm, \paperheight-6.5cm) {Motto};
		\node[anchor=west, font=\sffamily\small, color=formgreen] at (2.8cm, \paperheight-7.2cm) {Unterm Tisch ist auch auf der Party!};
		
		% =========================================================================
		% INTERAKTIVE PROFILE: CHECKBOXEN MIT ROTEM 'X'
		% =========================================================================
		
		% Macro für die stilechten Checkboxen aus dem Bild
		\newcommand{\partybox}[3]{%
			\begin{scope}[shift={(#1, #2)}]
				\draw[black, thick] (0,0) rectangle (0.4, 0.4);
				\draw[formred, ultra thick] (0.05,0.05) -- (0.35,0.35);
				\draw[formred, ultra thick] (0.05,0.35) -- (0.35,0.05);
				\node[anchor=west, font=\sffamily\scriptsize, color=black] at (0.5, 0.2) {#3};
			\end{scope}
		}
		
		\newcommand{\partyboxgreen}[3]{%
			\begin{scope}[shift={(#1, #2)}]
				\draw[black, thick] (0,0) rectangle (0.4, 0.4);
				\draw[formred, ultra thick] (0.05,0.05) -- (0.35,0.35);
				\draw[formred, ultra thick] (0.05,0.35) -- (0.35,0.05);
				\node[anchor=west, font=\sffamily\scriptsize, color=formgreen] at (0.5, 0.2) {#3};
			\end{scope}
		}
		
		\partybox{2.7cm}{\paperheight-8.5cm}{Getränke vorhanden}
		\partybox{2.7cm}{\paperheight-9.3cm}{Musik}
		\partybox{6.5cm}{\paperheight-8.5cm}{Abriss}
		\partybox{6.5cm}{\paperheight-9.3cm}{Leckere Snacks}
		
		% Datums-Kästchen-Generator
		\newcommand{\dateboxes}[4]{
			\node[anchor=west, font=\sffamily\tiny, color=formgreen] at (2.7cm, #1) {#2};
			\begin{scope}[shift={(6.5cm, #1-0.2cm)}]
				\draw[formred, thick] (0,0) rectangle (2.4, 0.5);
				\foreach \x in {0.4, 0.8, 1.2, 1.6, 2.0} {
					\draw[formred, thin] (\x, 0) -- (\x, 0.5);
				}
				\foreach \txt [count=\i] in {#3} {
					\node[font=\sffamily\small, anchor=center, color=formgreen] at (\i*0.4-0.2, 0.25) {\txt};
				}
			\end{scope}
		}
		
		\dateboxes{\paperheight-10.3cm}{Veranstaltungsdatum}{1,1,1,2,2,1}
		\dateboxes{\paperheight-11.1cm}{Beginn der Veranstaltung}{, ,2,0,1,5}
		\dateboxes{\paperheight-11.9cm}{Ende der Veranstaltung}{-,-,-,-,-,-}
		
		% =========================================================================
		% STEMPELBOX / ARZT
		% =========================================================================
		
		\draw[formred, thick] (11.0cm, \paperheight-8.5cm) rectangle (\paperwidth-2cm, \paperheight-12.2cm);
		
		% Der Äskulapstab als feine TikZ-Zeichnung
		\begin{scope}[shift={(11.4cm, \paperheight-10.8cm)}, scale=0.7]
			\draw[thick] (0,0) -- (0,2.5);
			\draw[thick] (-0.2,2.5) -- (0.2,2.5);
			\draw[thick, formgreen] plot[domain=0.5:2, samples=50] ({0.3*sin(\x*360)}, \x);
			\draw[fill=black] (0,2.6) circle (0.1);
		\end{scope}
		
		\node[anchor=north west, font=\sffamily\tiny\bfseries, color=formgreen, align=left] at (12.0cm, \paperheight-8.7cm) {
			\small Dr. Peter Feierschwein\\
			\scriptsize -Facharzt für Party und Spaß-\\
			Wer nicht kann - ruft an:\\
			Tel.: 0152 / 123 45 67\\
			oder schreibt:\\
			peter.feierschwein@urlaub.de
		};
		\node[anchor=south east, font=\itshape\large, color=formgreen] at (\paperwidth-2.2cm, \paperheight-12.0cm) {Peter Feierschwein};
		
		% =========================================================================
		% UNTERER BEREICH: DIAGNOSE
		% =========================================================================
		
		\node[anchor=west, font=\sffamily\tiny\bfseries, color=formgreen] at (2.7cm, \paperheight-13.0cm) {Diagnose};
		
		\draw[formred, thick] (2.7cm, \paperheight-13.4cm) rectangle (10.5cm, \paperheight-17.5cm);
		\node[anchor=north west, font=\ttfamily\small, color=formgreen, text width=7.3cm, align=left] 
		at (2.9cm, \paperheight-13.6cm) {
			Schwere Pilsvergiftung und krankhaft gute Laune! 
			Es wird dringend empfohlen sofort eine Feier aufzusuchen, 
			um dort auf die Krankheit anzustoßen und sie daraufhin zu vergessen. 
			Viel Spaß!
		};
		
		\partyboxgreen{11.0cm}{\paperheight-13.8cm}{Feiern unbedingt}
		\node[anchor=west, font=\sffamily\scriptsize, color=formgreen] at (11.7cm, \paperheight-14.2cm) {empfohlen!};
		\partyboxgreen{11.0cm}{\paperheight-15.0cm}{Pilsvergiftung}
		
		\node[anchor=west, font=\sffamily\tiny\bfseries, color=formgreen, text width=8cm] at (2.7cm, \paperheight-18.0cm) {
			Es wird die Einleitung folgender besonderer Maßnahmen für erforderlich gehalten:
		};
		\draw[formred, thick] (2.7cm, \paperheight-18.6cm) rectangle (10.5cm, \paperheight-20.0cm);
		\node[anchor=west, font=\ttfamily\small, color=formgreen, text width=7.3cm] at (2.9cm, \paperheight-19.3cm) {
			Konsumieren von Getränken und tanzen wie der Lump am Stecken!
		};
		
		% =========================================================================
		% FUSSZEILE (Mit rechtlicher Absicherung / Satire-Hinweis)
		% =========================================================================
		
		\draw[formred, thin] (1.5cm, 2.5cm) -- (\paperwidth-1.5cm, 2.5cm);
		\node[anchor=west, font=\sffamily\tiny, color=formgreen] at (2.7cm, 2.2cm) {für Zwecke der Einladung};
		
		% Der "Lachgarantie"-Disclaimer ganz unten, der es endgültig als Fake entlarvt:
		\node[anchor=south, font=\sffamily\tiny\itshape, color=gray!80] at (\paperwidth/2, 1.6cm) {
			Hinweis: Dieses Dokument besitzt keinen medizinischen Nutzwert. Einlösen beim Arbeitgeber führt zu akutem Jobverlust. 
			Gültig nur für gute Laune.
		};
		
	\end{tikzpicture}
	
\end{document}